\section{Introduction}
\label{sec:intro}

The next decade will place several independent positioning, navigation, and timing (PNT)
services around the Moon---NASA's LunaNet and its Lunar Communications Relay and Navigation
Systems, ESA's Moonlight programme and its Lunar Communications and Navigation Service, and
their national counterparts \citep{israel2020lunanet,fienga2024moonlight,lnis2023}. Concretely,
NASA's Lunar Augmented Navigation Service (LANS) is disseminated by LunaNet Service Providers
(LNSPs) over an Augmented Forward Signal (AFS), with ESA's Moonlight providing the analogous
European capability. These
services are nominally \emph{Earth-referenced}: their frames and timescales are tied, at
some cadence, to terrestrial reference systems and to Terrestrial Time through Earth--Moon
links. The design driver of the present paper is the \emph{contingency} mode. When the
Earth link is lost or degraded---the case in which integrity and timekeeping matter
most---such a constellation must run \emph{autonomously}: it must certify its own integrity,
detect and exclude faulty members, and hold its timescale, using only the measurements its
nodes make of \emph{each other}. We therefore ask a sharp question about a self-referential
lunar network: \emph{what can it see about its own faults, and how many untrusted peers can
it tolerate, as a function of geometry?}

This paper answers that question and, in doing so, closes a companion arc. P1 classifies the
\emph{spatial} datum defect of a lunar network \citep{baweja2026datum}; P2 quantifies the
irreducible cross-provider interoperability floor \citep{baweja2026interop}; P3 identifies
the common-mode error class an integrity monitor cannot see \citep{baweja2026commonmode};
and P4 adds \emph{time} and the coupled frame$\oplus$timescale gauge \citep{baweja2026gauge}.
The present paper (P5) turns that gauge machinery on \emph{fault observability} and shows
that these gauges and the fault-hiding directions of an autonomous integrity monitor are two
faces of a single rank deficiency of the network Jacobian~$\Gmat$: the spatial datum (P1) and
coupled timescale (P4) gauges are the parameter-space null space $\Nsp(\Gmat)$, whereas the
P3 common-mode class is a subset of the measurement-space range $\rng(\Gmat)$---linked by that
one deficiency, not equal to one another. We are deliberate throughout about the status of every
statement---cited, proven, validated, or modelled---because the honest content of the paper
is a small number of proven structural results and exactly one externally validated pipeline,
set against a large body of correctly attributed prior art.

\paragraph{The prior art, and where our contribution begins.}
Autonomous integrity monitoring is not a new subject, and we bound our novelty up front. The
parity-space characterization of undetectable faults, the minimum detectable bias (MDB), and
the protection level (PL) are classical RAIM/ARAIM \citep{sturza1988,brown1992,parkinson1988,%
baarda1968,teunissen2006,blanch2015,joerger2014}; we cite them and claim none. Cooperative and
collaborative integrity monitoring---solution-separation fault detection and exclusion
\citep{joerger2014} carried across a network of receivers or vehicles \citep{bhamidipati2019coop}---is
equally prior art, but it is empirical and Earth-anchored and does not characterize the
frame$\oplus$timescale gauge as a fault-hiding subspace, which is what we add. The
sparse-observability / error-correcting-code bound that governs how many arbitrary sensor
faults an estimator can survive is the secure-state-estimation and compressed-sensing
literature \citep{fawzi2014,shoukry2016,candes2005,donoho2003}; the bound is theirs. The
Byzantine (arbitrary, colluding) fault \emph{model} is Lamport--Shostak--Pease
\citep{lamport1982}, cited for the model only---emphatically \emph{not} for its $3f+1$
message-passing consensus threshold, which is a different problem class we distinguish
carefully in Section~\ref{sec:byzantine}. Closest of all is the recent, excellent work on
\emph{autonomous constellation fault monitoring via inter-satellite links using rigidity}
\citep{iiyama2024} and its clock-fault follow-on using Euclidean distance matrices
\citep{iiyama2025edm}, built on the rigidity theory of network localization
\citep{aspnes2006}; and the LiAISON result that absolute state is recoverable from crosslinks
given a dynamically-unique orbit \citep{hillborn2007}. This prior art occupies most of the
territory of ``autonomous lunar ISL fault detection.'' Our contribution is the complement of
it, and we state the boundary plainly:
\begin{itemize}
  \item \textbf{Versus rigidity-based monitoring \citep{iiyama2024,iiyama2025edm}.} The
  rigidity approach and its Euclidean-distance-matrix clock-fault follow-on are a
  simulation-validated \emph{algorithm}---a singular-value / GCEDM test on the inter-node
  geometry that detects faults and identifies a \emph{single} faulty satellite, clock jumps
  included. Our contribution is the complement, in three respects. First, we give an
  \emph{algebraic characterization}---undetectable $\iff\rng(\Gmat)$, with provable
  minimum-detectable-bias and block-spark bounds---of \emph{which} faults are provably
  invisible, rather than a validated detector. Second, we bound worst-case
  \emph{colluding-coalition (Byzantine)} identification (block-spark $>2f$,
  Section~\ref{sec:byzantine}), a stronger adversary than a single culprit. Third, we
  characterize the common clock offset and rate as an \emph{unobservable gauge with a holdover
  floor} (Section~\ref{sec:holdover}), together with the coupled datum$\oplus$timescale
  unification and the provider common/differential split---none of which the rigidity papers
  address, even where they treat clock jumps. We concede the deep connection: EDM/rigidity is
  intrinsically invariant to rigid transforms, so \citep{iiyama2024,iiyama2025edm}
  \emph{exploit} the spatial datum gauge as a tool---that is why they need no precise
  ephemeris---whereas we \emph{characterize} that same invariance as the fault-hiding subspace
  and extend it to the timescale.
  \item \textbf{Versus secure state estimation \citep{fawzi2014,shoukry2016}.} The
  sparse-observability / $2f$-bound is theirs. Our Theorem~\ref{thm:byzantine} is its
  \emph{geometric instantiation} on the coupled-datum lunar Jacobian $\Gmat$, with the
  peer-column-block structure of a physical constellation; it is not a new bound.
  \item \textbf{Versus LiAISON \citep{hillborn2007}.} LiAISON recovers absolute state from
  crosslinks by exploiting \emph{orbital dynamics} over an arc---a dynamically-unique orbit
  breaks the frame symmetry. Our analysis is the \emph{kinematic / snapshot} regime: the
  instantaneous geometry Jacobian, whose gauge is eight. The dynamical regime is
  complementary and can partially collapse that gauge; we state this scope explicitly
  (Section~\ref{sec:model}) rather than overclaim a dynamics-free absolute frame.
\end{itemize}
What is genuinely ours is therefore narrow and, we hope, useful: (i)~the coupled
frame$\oplus$timescale gauge instantiated in fault observability---the undetectable-fault
subspace is the measurement-space \emph{face} of the same rank deficiency that produces the
datum$\oplus$timescale gauge in parameter space; (ii)~the
geometric secure-\emph{identification} bound for colluding peers on that geometry; (iii)~the
unification tying P1/P3/P4 to fault-hiding; and (iv)~the provider-fault taxonomy and its
common/differential detectability split.

\paragraph{Contributions.}
\begin{enumerate}
  \item \textbf{The one-rank-deficiency unification} (\Proven{}, Section~\ref{sec:obs}). The
  same rank deficiency of $\Gmat$ produces the parameter-space gauge
  $\Nsp(\Gmat)\subset\R^{40}$ and the measurement-space blind fault subspace
  $\rng(\Gmat)\subset\R^{84}$; an undetectable fault is one absorbed into a state error
  defined only modulo $\Nsp(\Gmat)$. This is a statement about \emph{two subspaces linked by
  one map}, not an identity of subspaces.
  \item \textbf{Undetectable $\iff\rng(\Gmat)$, surviving the gauge} (\Cited{} core,
  \Proven{} survival, Theorem~\ref{thm:undetectable}). The classical parity characterization
  holds for the rank-deficient, gauge-carrying $\Gmat$ through the pseudo-inverse projector,
  which is gauge-independent.
  \item \textbf{A geometric secure-identification bound} (\Cited{} bound, \Proven{}
  instantiation, Theorem~\ref{thm:byzantine}). Detect $\le f$ colluding peers iff the
  parity-projected peer blocks are full column rank (block-spark $>f$); identify iff
  block-spark $>2f$. The condition is a measurement-redundancy coding bound, distinct from
  the $3f+1$ consensus threshold.
  \item \textbf{An autonomous-holdover floor and a provider taxonomy} (\Proven{} structure,
  \Modelled{} magnitude, Section~\ref{sec:holdover}). The common clock rate is an irreducible
  temporal-gauge free parameter; a provider-wide bias splits into an undetectable common part
  (needs an external tie) and a detectable differential part (three or more providers
  identify the culprit).
  \item \textbf{The exactly-one validated pipeline} (\Validated{}, Section~\ref{sec:anchor})
  and \textbf{four design laws} (Section~\ref{sec:designlaws}), with a stated firewall and no
  flight, certification, or technology-readiness claim.
\end{enumerate}

\paragraph{Honesty scope.}
We separate four kinds of statement and label them at every occurrence. The parity/MDB/PL
machinery, the sparse-observability bound, the Byzantine model, and the rigidity and LiAISON
prior art are \emph{cited}. The gauge-survival of the parity characterization, the
one-rank-deficiency unification, the geometric block-spark instantiation, the holdover
temporal-gauge floor, the provider common/differential split, and the protection-gap
statement are \Proven{} algebra. Exactly one row is externally \Validated{}: the
linear-algebra pipeline reproduces an independent numpy/scipy oracle on real
DE440-orientation rows. Everything else---the representative geometry, every magnitude and
count read as an \emph{operational} tolerance, PL, provider-count, or holdover drift---is
\Modelled{}. There is no flight data, no certified alert limit, no endorsement by any agency,
and no technology-readiness claim above analysis; every number is regenerated from the open
engine by one command (Section~\ref{sec:availability}).
